
%(BEGIN_QUESTION)
% Copyright 2012, Tony R. Kuphaldt, released under the Creative Commons Attribution License (v 1.0)
% This means you may do almost anything with this work of mine, so long as you give me proper credit

Jane har øreverk og bestemmer seg for å legge en varmeflaske mot øret for å lindre smerte. Som medisinstudent estimerer hun at det trengs et varmeinnslag på 19 000 kalorier.

\vskip 10pt

Anta at varmtvannet hennes holder 54 grader Celsius, og at hudtemperaturen er omtrent 31 grader Celsius. Beregn hvor mye vann hun må fylle i varmeflasken for å avgi den estimerte varmemengden når flasken kjøles fra starttemperatur til hudtemperatur.

\vskip 10pt

Avgjør også om beregnet verdi er et {\it høyt} eller {\it lavt} estimat, basert på forenklingene i beregningen.

\underbar{file i01011}
%(END_QUESTION)





%(BEGIN_ANSWER)

Vi skal finne vannmassen som trengs for å avgi 19 000 kalorier ved kjent temperatursenkning. Bruk spesifikk varme og løs for $m$:

$$Q = mc \Delta T$$

$$m = {Q \over c \Delta T}$$

$$m = {19000 \hbox{ cal} \over (1 \hbox{ cal/g-}^o\hbox{C}) (54^o\hbox{C} - 31^o\hbox{C})}$$

$$m = 826.1 \hbox{ gram}$$

1 liter vann er omtrent 1000 gram, så nødvendig mengde er 0.8261 liter.

\vskip 10pt

Dette er et {\it lavt} estimat, fordi ikke all varmen overføres til øret. Noe går til andre deler av hodet og til omgivelsesluften. Reelt behov er derfor høyere.

%(END_ANSWER)





%(BEGIN_NOTES)


%INDEX% Physics, heat and temperature: calorimetry problem 

%(END_NOTES)

